\section{Introduction}


\subsection{Subject}
 In October 2024, John J. Hopfield was awarded the Nobel Prize in Physics alongside Geoffrey Hinton for their
pioneering work on artificial neural networks \cite{nobel2024physics}. Their work laid the foundations for further development and understanding of neural networks, whether they are organic or artificial, and eventually gave rise to Artificial Intelligence (AI).\\
Hopfield worked on an artificial neural network originally invented by Shun'ichi Amari in 1972 \cite{Amari1972LearningPA}. He greatly contributed to the analysis, understanding and popularization of the network in his landmark paper \textit{Neural networks and physical systems with emergent collective computational abilities} published in 1982 \cite{Hopfield1982}, and the network was eventually named after him. The model, described in more details in subsection \ref{subsec:hopfield_model}, is similar to the Ising and the spin glass networks: it is composed of a set of spins -- called neurons in the case of the Hopfield network -- $S_i \in \{-1, 1\}$ interacting with one another. In the same way as the two preceding models, the Hopfield model is energy-based, that is it will tend to minimize its energy over time. What makes this model different from the other two is the way the interactions between neurons are defined: as for the Ising and spin glass models they remain static over time but depend on the set of patterns to be stored in the network, making the latter attractors, that is to say minima of energy, of the dynamics. \\
The Hopfield model provided the first simple yet powerful framework for understanding networks of interconnected neurons. It has become a reference in statistical mechanics on networks and its properties at equilibrium have been extensively studied and are now well known, whether it is its statistical mechanics \cite{Amit1985,AMIT1987,Amit1989,Sompolinsky1988} or its phase diagram and transitions \cite{Hertz1991}. However, much less is known about the dynamics that lead to those well-known equilibrium properties.

\subsection{Objectives}
The main objective of this internship is to efficiently simulate the Metropolis dynamics on the Hopfield network. To this end, we use a bitwise formalism that allows us to simulate several replicas of the system in parallel, increasing the efficiency of the exploration of the model's energy landscape. Our objectives in this internship are thus twofold: first, to validate the bitwise implementation and quantify its computational gain in yield over the standard Metropolis algorithm;
second, to leverage this implementation to study the dynamics of the Hopfield network. To validate the implementation, we begin with the Ising model, which serves as a natural benchmark due to its simplicity and well-characterized phase transition. Once this implementation has been proven sane and efficient on the Ising network, we increase the complexity by implementing it on the spin glass network, testing the gain in yield again.\\
Once the bitwise arithmetic has been validated on both models, we implement it on the Hopfield network in order to analyze its dynamics.

